\section{Open questions}\label{sec:open}

\begin{enumerate}[(1),leftmargin=*]
\item \emph{Five-row counts.} By Theorem~\ref{thm:periodic} the series
$\sum_{m\ge0}a_5(m)z^m$ is rational. By computer we know $a_5(5)=8$ and $a_5(m)$ for
$8\le m\le23$ (Section~\ref{sec:five}). Find this rational function, or a closed form for $a_5(m)$, and
the growth rate of $a_5(m)$.
\item \emph{Six rows below width $15$.} Theorem~A gives $f(6,m)\le m+8$ for
$8\le m\le14$. Stored SAT witnesses give $f(6,9)\ge16$ and $f(6,10)\ge17$
(\nolinkurl{jis/data/small_cells_sat_6x9_b16.json} and \nolinkurl{small_cells_sat_6x10_b17.json}),
so $f(6,9)$ is $16$ or $17$ and $f(6,10)$ is $17$ or $18$. For $11\le m\le14$ we know only
$m+5\le f(6,m)\le m+8$. What are these values?
\item \emph{Seven and more rows.} For $h\le6$ the row bound $m+q_h$ is attained at every
large width. Is this true for every $h$? By Theorem~\ref{thm:collar}, one valid blackout
of a $7\times m$ grid with $m+12$ points that passes the collar test with $W_7=10$ columns
would settle $h=7$. For $h\ge7$ we do not know whether $f(h,m)-m$ is eventually constant.
\item \emph{Squares.} We have $f(n,n)=2n$ for $n=5$, $6$, and $7$, the last two by
computer. What is the least $n$ with $f(n,n)>2n$? An alteration argument in the style of Alon
and Spencer~\cite[Section~3.2]{alonspencer2016} gives $f(n,n)\ge c\,n^{6/5}(1+\ln n)^{-2/5}$
for some $c>0$ and all large $n$. We do not include the proof here. It is in the file
\nolinkurl{jis/removed/appC_square.tex} of the repository~\cite{repo}. Neuwirth's

bound~\cite[Corollary~16]{neuwirth2001} for bipartite graphs of girth eight gives
$f(n,n)=O(n^{4/3})$. The true order of growth is open.
\item \emph{Proofs without a computer.} The $5\times5$ optimum is unique up to the
symmetries of the square, and the four-row generating function has a dominant pole of
order five. Both facts rest on computation. Is there a proof by hand of either one?
\end{enumerate}
